\chapter{凸関数}
    \section{諸定理}
        \subsection{一様ノルムの凸性}
            \begin{shadebox}
                $n\in\naturalNumbers$ とし，$S$ を $\realNumbers^n$ の部分集合とする。
                $f:S\to\realNumbers$ は有界であるとする。
                このとき一様ノルム $\norm{\cdot}_{\infty,S}: f\in S\mapsto\sup_{x\in S}\abs{f(x)}$ は凸である。
            \end{shadebox}
            \begin{proof}
                \quad\par
                $f_1,f_2:S\to\realNumbers$ は有界であるとし，$x\in S,t\in[0,1]$ とする。
                一様ノルムの定義の定義から $\norm{f_1}_{\infty,S}\geq\abs{f_1(x)},\norm{f_2}_{\infty,S}\geq\abs{f_2(x)}$ が成り立つので次式が成り立つ。
                \[ (1-t)\norm{f_1}_{\infty,S} + t\norm{f_2}_{\infty,S}\geq (1-t)\abs{f_1(x)} + t\abs{f_2(x)}\geq (1-t)\abs{f_1(x)} + t\abs{f_2(x)} = \abs{(1-t)f_1(x) + tf_2(x)} \]
                任意の $x\in S$ について上式が成り立つので $(1-t)\norm{f_1}_{\infty,S} + t\norm{f_2}_{\infty,S}\geq\norm{(1-t)f_1 + tf_2}_{\infty,S}$ が成り立つ。
            \end{proof}
        \subsection{\texorpdfstring{$f_1,f_2$}{f\_1,f\_2}が凸なら\texorpdfstring{$\max\{f_1,f_2\}$}{max(f\_1,f\_2)}は凸}
        \begin{proof}
            \quad\par
            $g(\bm{x}) \coloneq \max\{f_1(\bm{x}), f_2(\bm{x})\}$とする。
            $g$の定義域内の任意の$\bm{x}_1,\bm{x}_2$と任意の$\theta\in[0,1]$に対して
            \[ (1-\theta)g(\bm{x}_1) + \theta g(\bm{x}_2) = (1-\theta)\max\{f_1(\bm{x}_1), f_2(\bm{x}_1)\} + \theta\max\{f_1(\bm{x}_2), f_2(\bm{x}_2)\} \]
            であるが，maxの性質から
            \[ \max\{f_1(\bm{x}_1), f_2(\bm{x}_1)\} \geq f_1(\bm{x}_1), f_2(\bm{x}_1) \text{かつ} \max\{f_1(\bm{x}_2), f_2(\bm{x}_2)\} \geq f_1(\bm{x}_2), f_2(\bm{x}_2) \]
            であるので
            \[(1-\theta)g(\bm{x}_1) + \theta g(\bm{x}_2) \geq (1-\theta)f_1(\bm{x}_1) + \theta f_1(\bm{x}_2) \geq f_1[(1-\theta)\bm{x}_1 + \theta\bm{x}_2]\]
            かつ
            \[(1-\theta)g(\bm{x}_1) + \theta g(\bm{x}_2) \geq (1-\theta)f_2(\bm{x}_1) + \theta f_2(\bm{x}_2) \geq f_2[(1-\theta)\bm{x}_1 + \theta\bm{x}_2]\]
            である。このことから
            \[(1-\theta)g(\bm{x}_1) + \theta g(\bm{x}_2) \geq \max\{f_1[(1-\theta)\bm{x}_1 + \theta\bm{x}_2], f_2[(1-\theta)\bm{x}_1 + \theta\bm{x}_2]\} = g[(1-\theta)\bm{x}_1 + \theta\bm{x}_2]\]
        \end{proof}
        \subsection{Jensenの不等式}
            \begin{shadebox}
                $f(x)$を凸関数とし，$w_i\geq 0(i=1,\dotsc,n),\sum_{i=1}^n w_i=1$とするとき，次式が成り立つ。
                \[ f\left(\sum_{i=1}^n w_i x_i\right) \leq \sum_{i=1}^n w_i f(x_i) \]
            \end{shadebox}
            \begin{proof}
                \quad\par
                $n=1$の時は明らか。
                $n=2$のときは凸関数の定義より成り立つ。
                $n=n_0\in\naturalNumbers$のとき成り立つと仮定して$n=n_0+1$のとき成り立つことを示す。
                $w'_i\coloneq w_i/(1-w_{n_0+1})$とする(最初の$n_0$個の重みを規格化したもの)と$w'_i>0,w'_1+\dotsb+w'_{n_0}=1$だから，仮定より
                \begin{align*}
                    f\left(\sum_{i=1}^{n_0}w'_i x_i\right) &\leq \sum_{i=1}^{n_0}w'_i f(x_i) \\
                    \therefore\; f\left(\sum_{i=1}^{n_0}\frac{w_i}{1-w_{n_0+1}} x_i\right) &\leq \sum_{i=1}^{n_0}\frac{w_i}{1-w_{n_0+1}} f(x_i) \\
                    \therefore\; (1-w_{n_0+1})f\left(\sum_{i=1}^{n_0}\frac{w_i}{1-w_{n_0+1}} x_i\right) &\leq \sum_{i=1}^{n_0} w_i f(x_i) \\
                    \therefore\; w_{n_0+1}f(x_{n_0+1}) + (1-w_{n_0+1})f\left(X_{n_0}\right) &\leq \sum_{i=1}^{n_0+1} w_i f(x_i) \quad \left(X_{n_0} = \sum_{i=1}^{n_0}\frac{w_i}{1-w_{n_0+1}} x_i \text{とおいた}\right)
                \end{align*}
                $f$の凸性より
                \[ \text{左辺} \geq f(w_{n_0+1}x_{n_0+1} + (1-w_{n_0+1})X_{n_0}) = f\left(\sum_{i=1}^{n_0+1}w_i x_i\right)\]
            \end{proof}
        \subsection{重み付き相加,相乗平均の関係とYoungの不等式}
            \label{重み付き相加,相乗平均の関係とYoungの不等式}
            \begin{shadebox}
                $w_1,\dotsc,w_n\geq 0,w_1+\dotsb+w_n=1,x_1,\dotsc,x_n\geq 0$とするとき，次式が成り立つ。
                \[ \sum_{i=1}^n w_i x_i \geq \prod_{i=1}^n x_i^{w_i} \]
                この関係は\textbf{重み付き相加,相乗平均の関係}と呼ばれている。
            \end{shadebox}
            \begin{proof}
                \quad\par
                $x_i=0$である$i$が存在するときは明らかに成り立つ。
                以下では全ての$i$について$x_i>0$とする。
                対数関数の凹性を利用する。Jensenの不等式より
                \[ \log\sum_{i=1}^n w_i x_i \geq \sum_{i=1}^n w_i\log x_i \quad \therefore\; \log\sum_{i=1}^n w_i x_i \geq \log\prod_{i=1}^n x_i^{w_i} \quad \therefore\; \sum_{i=1}^n w_i x_i \geq \prod_{i=1}^n x_i^{w_i} \]
                最初の式より，等号成立条件は$x_1=\dotsb=x_n$である。
            \end{proof}
            この不等式には変形版があり，\textbf{Youngの不等式}と呼ばれている。
            一時的に上の定理の条件を少し強めて$w_1,\dotsc,w_n\textcolor{red}{>}0$とし，$y_i=x_i^{w_i}$とすると$x_i=y_i^{1/w_i}$となる。
            これらを重み付き相加,相乗平均の関係に代入して
            \[ \sum_{i=1}^n w_i y^{1/w_i} \geq \prod_{i=1}^n y_i \]
            を得る。$p_i=1/w_i$とすると次式を得る。
            \[ \sum_{i=1}^n \frac{{y_i}^{p_i}}{p_i} \geq \prod_{i=1}^n y_i \]
            改めて定理の形で述べると，次のようになる。\\
            \begin{shadebox}
                $p_1,\dotsc,p_n>1,1/p_1+\dotsb+1/p_n=1,x_1,\dotsc,x_n\geq 0$とするとき，次式が成り立つ。
                \[ \sum_{i=1}^n \frac{{x_i}^{p_i}}{p_i} \geq \prod_{i=1}^n x_i \]
            \end{shadebox}
        \subsection{重み付き幾何平均と重み付き調和平均の関係}
            \begin{shadebox}
                $n\in\naturalNumbers,x_i,w_i>0\;(i=1,\dots,n),\sum_{i=1}^n w_i = 1$ とするとき，$x_1,\dots,x_n$ の重み付き幾何平均は重み付き調和平均を下回ることはない。
                つまり次式が成り立つ。
                \[ \prod_{i=1}^n {x_i}^{w_i} \geq \frac{1}{\sum_{i=1}^n w_i/x_i} \tag{1} \]
            \end{shadebox}
            \begin{proof}
                \quad\par
                式 (1) は次式と等価である。
                \[ \sum_{i=1}^n w_i/x_i \geq \frac{1}{\prod_{i=1}^n {x_i}^{w_i}} \]
                ここで $\tilde{x}_i \coloneq x_i$ とすると上式は次式と等価である。
                \[ \sum_{i=1}^n w_i\tilde{x}_i \geq \prod_{i=1}^n {\tilde{x}_i}^{w_i} \]
                上式の成立は \ref{重み付き相加,相乗平均の関係とYoungの不等式} より明らかである。
            \end{proof}
        \subsection{2階導関数と凸性の関係}
            \begin{shadebox}
                $a\in\realNumbers$とし，$I$は$a$の開近傍であるとする。
                $f:I\to\realNumbers$は$I$上で$\mathrm{C}^2$級であるとする。
                次の3つの事項
                \begin{enumerate}
                    \item $\forall x\in I, f^{(2)}(x) > 0$
                    \item $\forall (x_1,x_2\in I,x_1\neq x_2),f(x_2) > f(x_1) + f^{(1)}(x_1)(x_2-x_1)$
                    \item $f$は$I$上で狭義に下に凸である。
                \end{enumerate}
                に対して$1\Rightarrow 2,2\iff 3$である。
                さらに$1\Leftarrow 2$は偽である。
            \end{shadebox}
            \begin{proof}
                \quad\par
                $1\Leftarrow 2$が偽であることは反例$f(x) = x^4$から直ちに示される。
                \par
                $1\Rightarrow 2$を示す。
                $x\in I,x\neq x_1$とする。
                $g(x) \coloneq f(x) - \bigl(f(x_1) - f^{(1)}(x_1)(x-x_1)\bigr)$とすると$g^{(1)}(x) = f^{(1)}(x) - f^{(1)}(x_1)$なので$g^{(1)}(x_1) = 0$である。
                さらに$g^{(2)}(x) = f^{(2)}(x) > 0$である。
                よって$g^{(1)}(x) < 0\;(x<x_1), g^{(1)}(x) > 0\;(x>x_1)$でるので$g(x_2) > 0$であり，2が成り立つ。
                \par
                $2\Rightarrow 3$を示す。
                $x_1,x_2\in I,x_1\neq x_2,0<\lambda<1,x_3 = (1-\lambda)x_1 + \lambda x_2$とすると次式が成り立つ。
                \[ f(x_1) > f(x_3) + f^{(1)}(x_3)(x_1-x_3),\quad f(x_2) > f(x_3) + f^{(1)}(x_3)(x_2-x_3) \]
                よって次式が成り立つ。
                \begin{align*}
                    (1-\lambda)f(x_1) + \lambda f(x_2) &> f(x_3) + (1-\lambda)f^{(1)}(x_3)(x_1-x_3) + \lambda f^{(1)}(x_3)(x_2-x_3) \\
                    &= f(x_3) + f^{(1)}(x_3)\left[(1-\lambda)(x_1-x_3) + \lambda(x_2-x_3)\right] \\
                    &= f(x_3) + f^{(1)}(x_3)\left[(1-\lambda)x_1 + \lambda x_2 - x_3\right] \\
                    &= f(x_3) + f^{(1)}(x_3)(x_3-x_3) = f(x_3)
                \end{align*}
                \par
                $3\Rightarrow 2$を示す。
                $x_1<x_2$の場合について示すが，逆の場合も同様にして示せる。
                $0<\lambda<1,x_3 = (1-\lambda)x_1 + \lambda x_2$とすると次式が成り立つ。
                \begin{align*}
                    \frac{f(x_2)-f(x_1)}{x_2-x_1} &= \frac{(1-\lambda)f(x_1) + \lambda f(x_2) - f(x_1)}{x_3-x_1} = \frac{(1-\lambda)f(x_1) + \lambda f(x_2) - f(x_1)}{\lambda(x_2-x_1)} \\
                    &> \frac{f(x_3) - f(x_1)}{\lambda(x_2-x_1)} = \frac{f(x_1 + \lambda(x_2-x_1)) - f(x_1)}{\lambda(x_2-x_1)} \tag{1} \\
                    &\to f^{(1)}(x_1)\; \text{as}\; \lambda\to 0 \\
                    \therefore \frac{f(x_2)-f(x_1)}{x_2-x_1} &\geq f^{(1)}(x_1) \quad \text{(真に大きいとまでは言えないことに注意)} \\
                    f(x_2) &\geq f(x_1) + f^{(1)}(x_1)(x_2-x_1)
                \end{align*}
                よって$x_4 \coloneq (x_1 + x_2)/2$とすると$f(x_4) \geq f(x_1) + f^{(1)}(x_1)(x_4-x_1)$となる。
                式(1)で$\lambda=1/2$とすると次式を得る。
                \begin{align*}
                    \frac{f(x_2)-f(x_1)}{x_2-x_1} &> \frac{f(x_4) - f(x_1)}{(x_2-x_1)/2} \\
                    f(x_2)-f(x_1) &> 2f(x_4) - 2f(x_1) \geq 2f^{(1)}(x_1)(x_4-x_1) = f^{(1)}(x_1)(x_2-x_1) \\
                    \therefore f(x_2) &> f(x_1) + f^{(1)}(x_1)(x_2-x_1)
                \end{align*}
            \end{proof}
    \section{Legendre変換}
        \subsection{十分滑らかな狭義凸関数に対してLegendre変換を2回施すと元の関数に戻る}
            \begin{proof}
                \quad\par
                $f$を開区間$I$で定義された$C^2$級の狭義凸関数とする。
                (よって1階導関数$f^{(1)}$は狭義単調増加だから逆関数が存在する。)
                $\fd$の値域を$I_2$とする。$f$のLegendre変換$f^*$は次のようになる。
                \[ f^*(p) = p\fd^{-1}(p) - f\left(\fd^{-1}(p)\right) \quad (p \in I_2) \]
                $f^*$のLegendre変換$f^{**}$を求めるためにまず${f^*}^{(1)}$の値域を調べる。
                \[ \fStarD(q)\;(q \in I_2) = \fd^{-1}(q) + q\frac{1}{f^{(2)}(q)} - \fd\left(\fd^{-1}(q)\right)\frac{1}{f^{(2)}(q)} = \fd^{-1}(q) \]
                より$\fStarD$の値域は$I$である。
                さらに，$\fd^{-1}$は狭義単調増加なので上式より$\fStarD$の逆関数が存在して$\fStarD^{-1} = \fd$である。
                $\fStarD^{-1}$の定義域は$I$である。
                以上より$f^{**}$は次のようになる。
                \begin{align*}
                    f^{**}(q)\; (q \in I) &= q\fStarD^{-1}(q) - f^*\left(\fStarD^{-1}(q)\right) \\
                    &= q\fd(q) - f^*(\fd) \\
                    &= q\fd(q) - \fd(q)\fd^{-1}\left(\fd(q)\right) + f\left(\fd^{-1}\left(\fd(q)\right)\right) \\
                    &= f(q)
                \end{align*}
            \end{proof}
            \let\fd\undefined
            \let\fStarD\undefined
    \section{諸注意}
        \subsection{非負の凸関数同士の積は凸関数とは限らない}
            反例として$x^2$と$(x+1)^2$の積がある。
            両者は単体では下に凸であるが，両者の積は$x=-1/2$で上に凸である。
